\chapter{Design and Specification of an MBSE framework for Nanosatellite EPS}
\label{ch:solucion}

The OOSEM activities that follow are concerned with specifying and designing the framework. They derive stakeholder needs, system requirements, measures, and architecture from direct communication with the stakeholders at ASCEND. Each level of abstraction (mission requirements, system requirements, logical architecture, physical architecture) specifies the next level of design, and produces test cases that support subsequent integration and verification \citep{friedenthal2014practical}.
At the request of the research group, certain stakeholder needs, system requirements, and use cases were ellicited from a literature review and then adjudicated to the stakeholders with their approval, given that they desire to develop capabilities in line with the academic state of the art.

% Based on feedback from the stakeholders, only the framework structure, interfaces, functions and use cases will be communicated via SysML v2 diagrams. They will be kept readable.
% Requirements, measures, verification activities, or more verbose content will be presented only using prose and tables.

\section{Analyze Stakeholder Needs}
\label{section:analyze-stakeholder-needs}
This activity characterizes the current operational domain, identifies stakeholder problems and their causes, and defines mission needs, measures of effectiveness, and the intended future domain. It establishes why the framework is needed and what successful use must achieve \citep{friedenthal2014practical}.

\subsection{As-is domain and causal analysis.}
The stakeholders are the members of the ASCEND research team involved in the EPS design process. The stakeholder roles identified are: Project Manager, Systems Engineer, EPS Engineer, Non-EPS Engineer, Testing Facility Technician. The Non-EPS Engineer is concerned with any other mission segment or subsystem apart from the EPS. This team structure is visible in the literature \cite{review-0529, review-1277, review-1155}, often working concurrently, and is generalized by \cite{review-0289} as "all systems architects and thematic experts".
Importantly, \cite{review-1438} acknowledges that in a small satellite project a single person may cover multiple roles, working under tight time and budget constraints.
% 0289 potential PA candidate, V&V example

% Problem statement from intro.tex
Evidence fragmentation resulting in low reproducibility of analyses and low capability to detect EPS design issues early is the initial problem that motivated this project.
However, the methodology OOSEM
proposes stakeholder analysis as an examination of the limitations of the existing EPS engineering systems and workflows from each stakeholder's perspective.
The concerns below incorporate feedback from ASCEND members and findings from the literature review.
Causal analysis is used to identify contributing factors and improvement opportunities, derive mission requirements, and establish measures of effectiveness that can quantify the value of a potential solution.

% Indentation is a temporary means to encode causality relations, is invisible in the compiled PDF so must be changed
% Effort is person-time

The Project Manager is concerned with:
\begin{itemize}
\item the high effort and cost to maintain verifiable requirements;
    \item the high effort to learn to write verifiable requirements;
    \item the high cost of requirements management tools used in the industry as a barrier to adoption, and;
\item the high effort to keep financial cost budgets and estimates up to date;
    \item the high effort to keep bills of materials up to date;
        \item the amount and complexity of candidate physical architectures compared;
        \item the amount of requirement and design changes;
\item the high cost and effort of redesign after hardware fabrication \citep[Secs.~4--5]{handpicked-0003};
    \item the late detection of thermal, electromagnetic compatibility, and power design issues \citep[Secs.~1 and~4]{handpicked-0003};
        \item the dependence of early verification on the availability of physical prototypes \citep[p.~3]{handpicked-0001};
\item the high cost of maintaining physical EPS replicas for mission troubleshooting \citep[Sec.~1]{review-0573};
\item the high effort and cost to adapt engineering tools to the team's mission scale and resources \citep[Sec.~5.1]{buruso2026modelbased};
    \item C-SE-2
\end{itemize}

The Systems Engineer is concerned with:
\begin{itemize}
\item the high effort to maintain traceability from requirements to design decisions;
    \item (C-SE-1-1) the high effort to learn how to use tools for requirements to design traceability;
        \item the high effort to learn how to use MBSE tools;
\item the high effort to communicate effectively requirement verification and validation evidence to stakeholders;
    \item the high effort to communicate effectively the requirements to design to verification evidence traces to stakeholders;
\item the high effort to maintain EPS requirements up to date when the designs of other subsystems change;
    \item the amount of other subsystem design changes;
        \item the amount of verification gates failed by other subsystem designs;
    \item C-EPS-1
    \item the high effort to define pass/fail criteria for EPS requirements;
        \item the high effort to translate the behavior and thermal and power characteristics of other subsystems into EPS requirement parameters;
            \item the high effort to determine the thermal and power characteristics of subsystem components;
            \item the large amount of possible representations of component behavior relevant to power consumption like duty cycles and operation modes;
            \item the large amount of possible methods to convert behavior and thermal and power characteristics into EPS requirement parameters, including power budgeting;
                \item the large amount of applicable standards;
                \item the large amount of methods published in the academic literature;
    \item the high effort to define verification activities for EPS requirements;
        \item the large amount of applicable standards;
        \item the large amount of potential verification activities published in the academic literature;
        \item the high effort to specify testing conditions and equipment, how they interact with the device under test, and the procedure to generate verification evidence in a format understandable by testing facility technicians;
            \item the large amount of variables that can be included in verification evidence to increase the translateability of evidence to an updated design in case of a failure;
\item the high effort to maintain power requirements for non-EPS subsystems up to date when the EPS design changes;
    \item the amount of EPS design changes;
        \item the amount of verification gates failed by the EPS designs;
\item the low reliability, availability, and maintainability (RAM) of CubeSat designs;
    \item the insufficient coverage of RAM variables in CubeSat and subsystem requirements;
    \item the high effort to verify requirements related to RAM variables;
        \item the high cost and effort to obtain evidence to verify requirements related to RAM variables;
            \item the proportion of requirements related to RAM variables that are verified via demonstration or testing instead of inspection or analysis;
            \item the high effort to perform RAM analyses on CubeSat and subsystem designs;
                \item the high effort to generate RAM analysis models of the CubeSat and subsystem designs;
\item the low capability to detect temporary mission energy shortages before hardware integration \citep[Sec.~VI]{review-1231};
    \item the use of average power budgets that omit interactions between attitude, solar generation, and changing load demand \citep[Secs.~I and~VI]{review-1231};
    \item the insufficient coverage of subsystem interactions and operating mode transitions in verification \citep[Sec.~8.1.7]{handpicked-0003}\citep[Sec.~3.1]{review-0573};
\item (C-SE-2) the high effort to maintain consistent information across requirements, simulation models, hardware, and verification evidence \citep[Secs.~4.4 and~5.1]{buruso2026modelbased};
    \item the high effort to reconcile incompatible modeling and simulation data formats \citep[Sec.~5.1]{buruso2026modelbased};
        \item the risk of semantic misalignment and information loss during manual data conversions \citep[Sec.~5.1]{buruso2026modelbased};
    \item the high effort to associate each verification result with the applicable design, software, model, and input data versions \citep[Sec.~4.4]{buruso2026modelbased};
    \item the high effort to carry measured behavior and revised model parameters back into design analyses \citep[Sec.~4.4]{buruso2026modelbased}\citep[p.~2]{handpicked-0001};
\item the high effort to establish that analysis results are credible verification evidence \citep[Sec.~6]{review-1251};
    \item C-EPS-2
    \item the high effort to define model validity ranges and acceptable prediction errors for the intended verification decisions \citep[Sec.~6]{review-1251};
    \item the uncertainty in predictions outside the tested loads, temperatures, and battery aging states \citep[pp.~5 and~8--9]{handpicked-0001};
\item (C-SE-3) the high effort to keep EPS hardware interfaces and flight software architecture consistent as the design changes \citep[Secs.~3--4]{review-0853};
    \item the difference between physical component groupings and the functional groupings used by flight software \citep[Sec.~3]{review-0853};
    \item the high effort to propagate EPS interface changes into software requirements and implementations \citep[Secs.~3--4]{review-0853};
        \item the late discovery of interface mismatches when hardware and software modules are developed separately \citep[Sec.~3]{review-0853};
    \item C-EPS-3
\item the high effort to coordinate EPS control, telemetry collection, fault recovery, and spacecraft safe mode \citep[Sec.~V]{citation-0136};
    \item the distribution of these responsibilities across separate flight software applications \citep[Sec.~V]{citation-0136};
    \item the high effort to specify how EPS commands and telemetry relate to spacecraft mode changes and recovery actions \citep[Sec.~V]{citation-0136};
\end{itemize}

The EPS Engineer is concerned with:
\begin{itemize}
\item (C-EPS-1) the high cost and effort to verify the EPS design after EPS requirements change;
    \item the high cost and effort to obtain evidence to verify EPS designs and compare candidate EPS physical architectures;
        \item the high effort to verify candidate COTS for EPS designs;
        \item the dependence of controller verification on the availability of complete physical prototypes \citep[p.~3]{handpicked-0001};
        \item the high effort to exercise controller changes under realistic electrical and timing constraints \citep[p.~3]{handpicked-0001}\citep[Sec.~3]{review-0573};
        \item the high effort to reproduce orbital changes in irradiance, temperature, and load demand on the test bench \citep[Secs.~3.2 and~4]{review-0573};
        \item the high cost and safety risk of obtaining evidence for rare or hazardous fault conditions \citep[p.~5]{handpicked-0001};
\item the high effort to maintain traceability from EPS requirements to EPS design decisions;
    \item (C-SE-1-1)
\item the high effort to translate the verification evidence of failed gates into an updated EPS design;
    \item the amount of variables relevant to EPS design included in verification evidence, which normally include only the inputs and outputs of the EPS as a black box;
    \item the high effort to maintain traceability from verification evidence to EPS design decisions;
\item the high effort to maintain EPS component requirements up to date when EPS designs change;
    \item the amount of EPS design changes;
        \item the amount of verification gates failed by the EPS designs;
    \item C-EPS-1
    \item the high effort to define pass/fail criteria for EPS component requirements;
        \item the high effort to determine the measurable variables of EPS components;
        \item the high effort to determine the allowable ranges of the relevant measurable variables of all EPS components;
    \item the high effort to define verification activities for EPS component requirements;
        \item the large amount of applicable standards;
        \item the large amount of potential verification activities published in the academic literature;
        \item the high effort to specify testing conditions and equipment, how they interact with the device under test, and the procedure to generate verification evidence in a format understandable by testing facility technicians;
            \item the large amount of variables that can be included in verification evidence to increase the translateability of evidence to an updated component choice or design in case of a failure;
        \item the high effort to specify inspection, analysis, and demonstration conditions, how they interact with the component under test, and the procedure to generate verification evidence;
            \item the large amount of variables that can be included in verification evidence to increase the translateability of evidence to an updated component choice or design in case of a failure;
            \item compatibility of specialized EPS engineering tools, like CAD and electrical simulators, with inspection, demonstration, and analysis procedures;
\item (C-EPS-2) the high effort to keep EPS behavioral models representative of the physical design \citep[p.~2]{handpicked-0001}\citep[Sec.~6]{review-1251};
    \item the high effort to calibrate nonlinear component models against experimental data \citep[Sec.~4.1]{review-1251};
        \item the amount of interacting model parameters and calibration measurements \citep[Sec.~4.1]{review-1251}\citep[p.~8]{handpicked-0001};
    \item the divergence of model predictions from hardware behavior as operating conditions and component health change \citep[p.~2]{handpicked-0001};
        \item the delay between collecting measurements and updating model parameters \citep[p.~2]{handpicked-0001};
    \item the computational cost of coupled electrical, thermal, and degradation models \citep[pp.~2--3]{handpicked-0001};
        \item the conflict between model fidelity and the execution time available for interaction with hardware \citep[pp.~2--3]{handpicked-0001}\citep[Sec.~3]{review-0573};
    \item C-SE-2
\item the high uncertainty in battery energy availability and remaining lifetime under mission loads \citep[Sec.~3]{handpicked-0003}\citep[pp.~1--3]{handpicked-0001};
    \item the limited visibility of internal battery states, including state of charge and state of health, from directly measured signals \citep[p.~1]{handpicked-0001};
    \item the omission of long term degradation mechanisms from simplified battery models \citep[pp.~1--3]{handpicked-0001};
    \item the coupled effects of load profiles, temperature, and depth of discharge on battery aging \citep[Sec.~3]{handpicked-0003}\citep[p.~3]{handpicked-0001};
\item the low capability to detect EPS anomalies before protection thresholds are exceeded \citep[Sec.~5.1]{review-1251};
    \item the high effort to determine the expected EPS response from operating states and telemetry spread across subsystem views \citep[Sec.~5.1]{review-1251};
    \item the high effort to distinguish component faults from normal degradation and measurement error \citep[Sec.~6]{review-1251};
\item the high effort to verify that battery charging and power management decisions are safe and effective \citep[pp.~2 and~6]{handpicked-0001};
    \item the gap between accurate battery behavior predictions and evidence that a control policy improves hardware performance \citep[p.~2]{handpicked-0001};
\item (C-EPS-3) the high effort to maintain consistent EPS command, telemetry, event, and parameter definitions across models, flight software, and ground software \citep[Secs.~3--4]{review-0853};
    \item the high effort to separately implement command, telemetry, event, and parameter definitions in the model export workflow described by \citet[Sec.~3]{review-0853};
    \item the high effort to translate modeled EPS behavior into mission-specific software functions \citep[Sec.~3]{review-0853};
\item the high effort to specify and verify EPS command execution conditions and expected effects \citep[Secs.~4--5]{review-0853};
    \item the high effort to translate EPS behavior into command sequences and confirm their functional flow with the flight software team \citep[Sec.~5]{review-0853};
\item the high effort to communicate EPS designs during technical reviews;
    \item the high effort to generate readable, accurate, and informative graphical representations of EPS designs;
\end{itemize}

The Non-EPS Engineer is concerned with:
\begin{itemize}
\item the high cost and effort to verify a subsystem design;
    \item the high cost and effort to verify the power requirements of a subsystem design;
        \item the high effort to reproduce the subsystem's dynamic current demand during mission operations \citep[Sec.~8.1.7]{handpicked-0003};
        \item the high effort to verify subsystem behavior during EPS supply changes and protection events \citep[Secs.~8.1.6--8.1.7]{handpicked-0003};
    \item C-SE-2
\item the high effort to integrate EPS services with flight software frameworks such as cFS and F-Prime \citep[Secs.~3--5]{review-0853}\citep[Sec.~III.B]{citation-0136};
    \item C-SE-3
    \item C-EPS-3
    \item the high effort to map EPS interfaces onto framework-specific communication mechanisms \citep[Sec.~3]{review-0853};
        \item the difference between typed component connections in F-Prime and publish-subscribe communication in cFS \citep[Sec.~3]{review-0853};
    \item the high effort to configure and deploy the selected framework on the OBC processor and operating system \citep[Sec.~III.B]{citation-0136};
        \item the high effort to learn framework configuration and integration tools \citep[Sec.~III.B]{citation-0136};
        \item the uncertainty over compatibility between the selected framework version and the target operating system \citep[Sec.~III.B]{citation-0136};
    \item the high effort to verify processor, memory, timing, and power budgets for EPS flight software functions \citep[Secs.~III.A.3, IV.A.1 and~VI.B--C]{citation-0136};
        \item the processing and memory limits of the OBC \citep[Secs.~IV.A.1 and~VI.C]{citation-0136};
        \item the interaction between operating system capabilities, processor performance, and power consumption \citep[Sec.~III.A.3]{citation-0136};
\end{itemize}

The Testing Facility Technician is concerned with:
\begin{itemize}
\item the high effort to assess the feasibility and safety of a proposed verification activity that takes place in the testing facilities;
    \item (C-TFT-1-1) the small amount of information about the procedure to generate verification evidence tipically included in a verification activity proposal;
    \item the small amount of safety-relevant information about the device under test tipically included in a verification activity proposal;
    \item the small amount of information about the interfaces between testing equipment and the device under test tipically included in a verification activity proposal;
    \item the risk of damaging batteries or power electronics while exercising fault conditions \citep[p.~5]{handpicked-0001}\citep[Sec.~8.1.6]{handpicked-0003};
        \item the high effort to establish safe fault injection limits before applying them to hardware \citep[Sec.~8.1.6]{handpicked-0003};
\item the high effort to generate verification evidence;
    \item the amount of variables included in the verification evidence;
    \item the complexity of the evidence generation procedures;
    \item the high effort to synchronize applied stimuli, hardware responses, and telemetry timestamps \citep[Sec.~8.1.7]{handpicked-0003}\citep[Sec.~6]{review-1251};
        \item the sampling, calculation, and output delays between simulated and physical equipment \citep[Sec.~3]{review-0573};
    \item the high effort to check the completeness and validity of recorded telemetry \citep[Sec.~8.1.7]{handpicked-0003}\citep[Sec.~4.2]{review-1251};
        \item the presence of invalid sensor values when components are inactive \citep[Sec.~4.2]{review-1251};
\end{itemize}

ASCEND is currently associated to a Rocketry Team and an Avionics Team of a student club.
The Rocketry Team operates on requirements given by Rocketry Competition Organizers.
The requirements are stated in terms of the target maximum height or maximum acceleration of the rocket.
One competition imposed the requirement of launching and returning to ground an egg without breaking it, which can only be validated via demonstration.
The Team Captains convert the requirements of a competition to atomic tasks.
Tasks are performed by any club member and approved by senior members.
The club performs the PDR, CDR, and FRR review gates.

There is a large emphasis on technology reuse at the club, so most tasks are formulated as atomic edits to existing designs.
Furthermore, most of the tasks for a particular competition involve development activities like PCB manufacturing, soldering, and assembly.
One observed task involved adding a COTS to a PCB, which was specified in the task title. The task was broken down into subtasks, one of which required to connect a capacitor to the COTS, possibly prescribed by its datasheet.

The competition requirements include the radio communications technologies allowed for rocket launch commands and telemetry.
Thus, radio transceivers are some of the first COTS chosen for rocket designs.
The egg mission also introduced a payload consisting of 1U unit with a mechanically damping material to protect the egg, and heat pumps to maintain its temperature within a range inspired by an incubator machine.

Multiple lithium-ion cells were used to power the latest revision of the egg incubator payload.
To size them, a power budget was generated by consulting the datasheets of the already chosen components and extracting "on" and "off" power consumption values, then estimating the duty cycle of each component for one launch.
The duty cycle of telemetry components is estimated from the telemetry frequency requirements of the competition. For the rest of the components, it is conservatively assumed that they are always on.

A PCB design was observed consisting of multiple boards connected perpendicularly via PCI to a backplane that was the size of a 1U face.
Status indicator LEDs were observed for 3.3 and 5 volt buses.
The power source of the payload is verified experimentally.
The payload is fully assembled and set to provide telemetry at a realistic frequency while powered by the battery.
It is then inserted into a temperature controlled environment like a refrigerator so the heat pumps can be observed under representative operation conditions, and left to operate until the battery power runs out.
If the time from experiment start to power depletion is greater than the mission period, the power requirements are taken to be verified.
In practice, the battery life is several times larger than the mission period.

The demonstratable requirement to protect an egg has been related to the measurable variable of acceleration in the egg chamber. An IMU was placed inside this chamber, which is mechanically damped, and another in a rigid point of the payload chassis.
The difference in acceleration between the two IMUs during realistic launches is analyzed to iterative on the damping design.
Launches are recreated by manually dropping the payload from a representative height.
The drop data, and naturally an inspection of the egg, is used to characterize the damping as "too soft" or "too hard".

Maximum rocket height and maximum rocket acceleration requirements are verified at the design phase using the OpenRocket software.
This software works with abstract representations of rocket components (nose, stages, etc.) can be parametrized, combined, and then simulated.
Rocket height and acceleration requirements are taken to be verified based on the results of the simulatins, and only then does development start.

The club reports that it takes about 2 months to onboard new members to the point where they can complete tasks.
Separately, ASCEND researchers have spent 9 months on one iteration of an EPS design following a more traditional systems engineering approach, starting from requirements and currently advenced up to high-level design, with a power budget completed and battery and panel options compared but no component selections yet.

\section{Analyze System Requirements}
This activity specifies the framework as a black box. It derives requirements from stakeholder needs and defines mission scenarios, system context, interfaces, behavior, performance, states, failures, and design constraints \cite{friedenthal2014practical}.

\section{Define Logical Architecture}
% Must be inspired at least by El Hachimi et al. (2025) and the CubeSat Handbook Chapter 1.

\section{Synthesize Candidate Physical Architectures}
This activity maps logical components to realizable hardware, software, persistent data, and procedure components. It first distributes logical responsibilities among nodes, then allocates them to physical components and records the rationale for important alternatives \citep[Sec.~17.3.5]{friedenthal2014practical}.

\section{Support System Integration and Verification}
This activity supports system integration and verification by defining test cases, verification procedures, environments, and criteria. Executable and analytical models support early, virtual integration and verification before physical implementation \cite{friedenthal2014practical}.

% proposal = anteproyecto
A project proposal was written before a complete understanding of the stakeholders' needs was achieved. Said proposal defined objectives, deliverables, indicators, and deadlines that form the research contribution baseline, without specifying the framework itself.
The assessment of the objective deliverables against their indicators and deadlines will be carried out separately from the framework integration and verification that happens within OOSEM (Section \ref{section:integrate-and-verify-system}).
% I will add the deliverables, indicators, and deadlines when the Solution section is complete. I will also add a subsection to the Results section for the indicator assessments.
